\section{Preliminary Results}\label{sec:evaluation}
\subsection{Workloads and experimental setting}
The experiments examine whether joint design improves executable workloads, which software--hardware interactions account for the gains, and how revising the formulation helps search. Two platforms provide complementary evidence. A programmable architecture benchmark exposes compute arrays, register and shared storage, DMA concurrency, data layout, and attention organization. Experiments on the open-source Ventus GPU architecture compile OpenCL kernels to instruction streams and explore parameterized compute, cache, register, and memory-service resources. Each study compares candidates against its own fixed baseline under common resource limits.

The programmable benchmark uses five complete Transformer scenarios: two prefill workloads and three decode workloads with different batch sizes and context lengths. All five programs share one hardware configuration under a 100-unit modeled area budget and the benchmark's power limits. Private-seed checks validate numerical outputs and resource legality. The Ventus experiments use Qwen2.5-0.5B with 24 layers and a 1.11\,mm$^2$ modeled area budget. P128 and P512 denote prefill lengths; D128 and D512 average 16 decode steps starting at those context lengths. Selected designs complete instruction-level timing simulation of every layer and decode step, including packing, transfers, and the language-model head.

For Ventus scenario $w$, $r_w=L_w/L_w^{\mathrm{ref}}$ measures normalized cycles. Each task minimizes the worst ratio over its prefill, decode, or combined scenario set. The programmable benchmark instead reports a weighted geometric-mean speedup. Results are simulated cycles and modeled resources. Ventus executes control and address paths for timing, with separate local numerical checks; the programmable benchmark checks full workload outputs. The following experiments establish design mechanisms and executable candidates for the proposed agent loop. Appendix~\ref{app:evidence} records the normalization and evidence definitions.

\subsection{Substantial gains from joint architecture and program design}\label{sec:joint-results}
The programmable benchmark reaches a weighted geometric-mean speedup of $7.02\times$ over its supplied baseline. Individual speedups are $8.37\times$ and $13.00\times$ for the two prefill cases, and $4.02\times$, $5.22\times$, and $5.01\times$ for decode (Figure~\ref{fig:controlled}a). The same selected hardware serves all five scenarios, so the aggregate includes the tradeoff between long attention, weight reuse, and decode concurrency. Compute-array shape, register capacity, data movement, and program organization are chosen together within the benchmark constraints.

On Ventus, a short exploratory campaign spanning three design iterations achieves $2.93\times$ prefill and $3.58\times$ decode speedups. The prefill and combined objectives select a design of estimated area 1.0812\,mm$^2$, with P128/P512 ratios of 0.3417/0.3308. The decode-oriented design occupies 1.0848\,mm$^2$ and reaches D128/D512 ratios of 0.2791/0.2578. The headline task ratios are determined by the 128-token scenarios and confirmed by complete instruction runs.

Controlled comparisons in the second iteration isolate each side's contribution. With the common slice evaluator, software-only optimization reaches prefill/decode objectives of 0.8581/0.9015, while hardware-only optimization reaches 0.4409/0.4325. Allowing both to change reaches 0.3656/0.2789 (Figure~\ref{fig:controlled}b). Software-first selection ties joint selection in this evaluated menu. The comparison establishes the value of changing both sides; the search policy's additional value depends on the interactions and alternatives represented in its domain.

The software gains themselves depend on hardware. Across the reported comparisons, software changes reduce cycles by approximately 10--14\% on the baseline machine. On the fixed optimized machine, paired P512 and D512 comparisons reduce cycles by approximately 18\% and 41\%. Column ordering, packing, and decode specialization therefore have greater value when the machine can support their reuse and memory-service demands.

\begin{figure}[tbp]
\centering
\begin{tikzpicture}
\begin{axis}[width=6.7cm,height=5.7cm,ybar,bar width=14pt,ymin=0,ymax=15,ytick={0,5,10,15},ylabel={Speedup over baseline},symbolic x coords={W-P,A-P,W-D1,W-D16,W-D4L},xtick=data,x tick label style={rotate=25,anchor=east},enlarge x limits=.15,title={(a) Programmable architecture benchmark},title style={font=\small},tick label style={font=\scriptsize},label style={font=\footnotesize},nodes near coords,nodes near coords style={font=\scriptsize},every node near coord/.append style={/pgf/number format/fixed,/pgf/number format/precision=2},ymajorgrids,grid style={gray!15}]
\addplot[fill=blueedge!70,draw=blueedge] table[x=label,y=speedup,col sep=comma] {figures/data/phase2.csv};
\end{axis}
\begin{axis}[at={(8.4cm,0)},anchor=south west,width=7.1cm,height=5.7cm,xbar,bar width=7pt,xmin=0,xmax=1.06,xtick={0,.25,.5,.75,1},xlabel={Worst normalized cycles},symbolic y coords={Joint,SW then HW,HW only,SW only,Baseline},ytick=data,enlarge y limits=.15,title={(b) Ventus Qwen2.5-0.5B controls},title style={font=\small},tick label style={font=\scriptsize},label style={font=\footnotesize},legend style={font=\scriptsize,at={(.5,-.32)},anchor=north,legend columns=2,draw=none},xmajorgrids,grid style={gray!15}]
\addplot[fill=blueedge!70,draw=blueedge] table[x=prefill,y=method,col sep=comma] {figures/data/ventus-controls-05b.csv};
\addplot[fill=greenedge!65,draw=greenedge] table[x=decode,y=method,col sep=comma] {figures/data/ventus-controls-05b.csv};
\legend{Prefill,Decode}
\end{axis}
\end{tikzpicture}
\caption{Preliminary co-design results. (a) One hardware configuration serves five complete Transformer scenarios and attains a $7.02\times$ weighted geometric-mean speedup. Higher is better. (b) Second-iteration Ventus controls use the same slice evaluator and area budget; varying both software and hardware improves over changing either alone. Lower is better. Complete runs confirm the selected task results; the third iteration further improves prefill to $2.93\times$.}\label{fig:controlled}
\end{figure}


\subsection{Coupled choices require an evolving formulation}\label{sec:tiling-results}
The performance gains expose concrete dependencies between program structure and resource allocation. Column-order GEMM places blocks that consume a common weight strip close together. Cache capacity then becomes useful where row-order sweeps had shown little benefit. Hardware search reallocates area toward outstanding memory requests and load/store queues. A decode-specific kernel loads only the nonzero fragment register for a single-row input, releasing load/store capacity without changing the matrix computation.

Larger tiles reveal a stronger interaction. Traffic analysis of a prefill projection finds 57\,MB of DRAM reads for 8.7\,MB of weights because row tiles reload the same data. Introducing a $64\times64$ GEMM kernel reduces traffic by roughly one third, but requires 173 vector registers per thread. On the earlier VGPR-512 configuration, normalized P128 cycles worsen from 0.362 to 0.438. A matching configuration with two multiprocessors, VGPR 1024, tensor shape $(4,8,4)$, and revised caches makes the tile useful within the area budget.

Hardware changes also reverse an earlier software decision. Packed QK is unfavorable on the baseline and the earlier three-multiprocessor design, yet improves P512 from approximately 0.377 to 0.334 on the new machine. A subsequent L2 adjustment produces the selected 0.3417 prefill objective. Figure~\ref{fig:trajectory} shows the accepted complete-simulation trajectory. Eliminating packed QK permanently after its first evaluation would remove a useful implementation from this later search space.

\begin{figure}[tbp]
\centering
\begin{tikzpicture}
\begin{axis}[width=.89\linewidth,height=5.6cm,ymin=0,ymax=.45,ytick={0,.1,.2,.3,.4},ylabel={Prefill worst normalized cycles},symbolic x coords={Earlier,N3,N4},xtick=data,xticklabels={Earlier accepted design,New tile + hardware,L2 variant},enlarge x limits=.2,ybar,bar width=35pt,nodes near coords,point meta=y,nodes near coords style={font=\footnotesize},every node near coord/.append style={/pgf/number format/fixed,/pgf/number format/precision=4},tick label style={font=\footnotesize},label style={font=\footnotesize},ymajorgrids,grid style={gray!15}]
\addplot[fill=blueedge!65,draw=blueedge] table[x=design,y=prefill,col sep=comma] {figures/data/trajectory.csv};
\end{axis}
\end{tikzpicture}
\caption{Accepted-design trajectory from complete instruction simulations. Introducing the larger-tile kernel and a matching register/tensor/cache configuration expands the available combination; an L2 variation then changes the selected hardware. The values are rounded as reported. This trajectory illustrates design evolution, rather than a controlled ablation of individual changes.}\label{fig:trajectory}
\end{figure}


The formulation evolves alongside these implementation choices. An initial search stalls near a normalized prefill estimate of 0.433. Factorial experiments expose interactions among load/store concurrency, L2 miss-status entries, and write buffering; adding pairwise terms to the MIP helps reach approximately 0.362, a further 16\% reduction. The larger tile then requires expanding the candidate family. The existing surrogate's penalty for poorly sampled choices keeps it from proposing the new combination, which targeted simulation subsequently finds and refines.

These observations connect representation construction to optimization. Program transformations change the resources worth buying, while hardware changes which transformations should remain available. An agent can respond by adding candidates, releasing a coupled neighborhood, and revising the interactions encoded in the formulation. Solving a fixed model more accurately alone cannot expose a kernel or hardware combination missing from that model.

\subsection{Joint selection improves the storage--latency tradeoff}\label{sec:storage-results}
A separate Ventus experiment minimizes allocated storage for one Qwen KV-append step and all 14 attention heads of a layer. It combines ten area-feasible hardware candidates with software choices for K/V precision, quantization group size, grouped-query organization, and scratch reuse. The analysis covers 2,240 completed grid points. Storage accounting includes quantization scales, scratch buffers, alignment, and private stacks. Independent probe seeds check local attention-output error against the BF16 path.

At normalized latency limits around 0.18--0.20, joint selection saves approximately 16\% or 17\% of storage over software-only optimization under 2\% or 5\% local error thresholds, respectively. At tighter limits around 0.15--0.17, software-only selection becomes infeasible while joint designs remain feasible. Faster hardware preserves the option to use a more compact format under the latency constraint. Hardware-first selection matches the joint solution in this grid, giving a useful strong baseline for future search comparisons.

The result extends the co-design objective beyond speed. Hardware configuration, representation precision, and temporary-storage organization jointly determine which designs satisfy a workload contract. The local error thresholds apply to this attention subproblem; whole-model quality is a separate deployment constraint.

\subsection{Execution feedback identifies missing model structure}\label{sec:model-results}
An independently constructed Ventus search provides a complementary diagnostic. Nine MIP solves close the gap in their approximate models and predict a worst-scenario ratio near 0.3333. Complete instruction simulation instead yields objectives between 0.7913 and 0.8987. Hardware-only selection matches or outperforms joint selection in that candidate set. Transferring the selected Qwen configurations to GPT-2 and Pythia also reverses the gains, with additional time concentrated in the FFN down projection and substantially increased memory reads.

This discrepancy identifies a concrete role for feedback. Solver optimality describes the current formulation; execution tests whether its traffic, reuse, and concurrency assumptions rank actual candidates correctly. The successful iterative study shows how adding such interactions and revisiting implementations improves the design. A progressive agent should preserve stronger evaluated candidates while changing the formulation responsible for a bad prediction.

Together, the results demonstrate substantial preliminary gains, hardware-dependent software preferences, and a storage-feasibility benefit from joint selection. They also provide specific actions for the proposed agent: expand an implementation family, revise coupling terms, and reopen a previously unfavorable branch. Evaluating the complete agent policy will compare these actions under common candidate access and total exploration budgets, measuring both resulting design quality and the cost of obtaining it.
